\subsection{一般注记}

设 A 是拓扑空间 E 的一个子集。为使 A 中点的一个序列 \((x_{n})_{n\in\mathbf{N}}\) 以 E 的一个点 x 为极限点，必须且只需下列条件成立 (GT, I, § 7, No.~3)：

\begin{enumerate}
\def\labelenumi{(\Alph{enumi})}
\tightlist
\item
  对每个整数 \(m \geqslant 0\) 以及每个邻域 \(\mathbf{U}\)（\(x\) 的一个邻域），存在一个整数 \(n \geqslant m\) 使得 \(x_{n} \in \mathbf{U}\)。
\end{enumerate}

形如 \((y_{k})_{k\in\mathbb{N}}\)（其中 \(y_{k}=x_{n_{k}}\)，\((n_{k})_{k\in\mathbb{N}}\) 是严格递增的正整数序列）的序列称为序列 \((x_{n})_{n\in\mathbb{N}}\) 的一个子序列。若序列 \((x_{n})_{n\in\mathbb{N}}\) 存在一个收敛到 x 的子序列，则 x 是 \((x_{n})\) 的一个极限点；反之，若 x 具有可数的基本邻域系，且 x 是序列 \((x_{n})\) 的极限点，则存在 \((x_{n})\) 的一个收敛到 x 的子序列。

根据 GT, IX, § 2, No.~9 的推论，我们断定，当 E 可度量化时，下列条件等价：

\begin{enumerate}
\def\labelenumi{\alph{enumi})}
\item
  集 A 在 E 中相对紧；
\item
  A 中点的每个无穷序列在 E 中有一个极限点；
\item
  从 A 中点的每个无穷序列，可以抽取出一个收敛到 E 中某点的序列。
\end{enumerate}

在本节中，我们将把这个判别法推广到某些不可度量化的拓扑向量空间。下面的命题使我们在许多情形下能把紧集的研究归结为弱紧集的研究。

命题 1. --- 设 E 是一个 Hausdorff 局部凸空间，A 是 E 的一个子集。设 \(E_{\sigma}\) 表示赋以弱化拓扑的空间 E。

\begin{enumerate}
\def\labelenumi{\alph{enumi})}
\item
  若 \(\mathbf{A}\) 中点的每个无穷序列在 \(\mathbf{E}\) 中有一个极限点，则 \(\mathbf{A}\) 在 \(\mathbf{E}\) 中预紧。
\item
  为使 A 在 E 中相对紧，必须且只需它在 E 中预紧且在 \(E_{\sigma}\) 中相对紧。
\end{enumerate}

我们用反证法证明 \(a)\)。若 \(A\) 不是预紧的，则由 GT, II, § 3, No.~7 的定理 3，可知存在一个对称凸邻域 \(V\)（\(E\) 中 0 的邻域），使得 \(A\) 不能被 \(V\) 的有限个平移所覆盖。

换句话说，若 \(x_{0}, x_{1}, \ldots, x_{n-1}\) 是 A 的点，则 \(A \not\subset \bigcup_{0 \leqslant i < n} (x_{i} + V)\)，于是存在 A 的一个点 \(x_{n}\) 使得 \(x_{n} - x_{i} \notin V\) 对 \(0 \leqslant i < n\) 成立。然后，对整数 n 用归纳法，可以构造 A 中点的一个无穷序列 \((x_{n})_{n \in \mathbb{N}}\)，使得 \(x_{n} - x_{m} \notin V\)（当 n \textgreater{} m 时）；由于 V 是对称的，我们还有 \(x_{m} - x_{n} \notin V\)（对 \(m \neq n\)），从而诸集 \(x_{n} + \frac{1}{2}V\) 两两不相交。对 E 中的每个点 x，至多存在一个整数 \(n \geqslant 0\) 使得 \(x_{n} \in x + \frac{1}{2}V\)，因而序列 \((x_{n})_{n \in \mathbb{N}}\) 没有任何极限点。这就证明了 a)。

现在假设 A 在 E 中预紧，并且包含在 \(E_{\sigma}\) 的一个紧子集 B 中。那么 B 在 \(E_{\sigma}\) 中完备，从而在 E 中也完备 (IV, 第 5 页, 注记 2)。我们有 \(\overline{A} \subset B\)，故 A 在 E 中相对紧。逆命题是显然的，b) 随之得证。

